% Auto-generated from raw.txt
% Usage:
%   \vocabentry{word}{/ipa/ (pos)}{definition}{example}{note}

\vocabentry{Recruitment}
{/rɪˈkruːtmənt/ (n)}
{The action of finding new people to join an organization.}
{The company has tightened its recruitment process to ensure they hire the best talent.}
{Đây là một danh từ dùng để chỉ quá trình tuyển dụng nhân sự mới.}

\vocabentry{Occupation}
{/ˌɑːkjuˈpeɪʃn/ (n)}
{A job or profession.}
{Please state your name, age, and occupation on the form provided.}
{Đây là một danh từ trang trọng dùng để chỉ nghề nghiệp hoặc công việc mà một người đang làm.}

\vocabentry{Redundancy}
{/rɪˈdʌndənsi/ (n)}
{The state of being no longer in employment because there is no more work available.}
{The merger of the two banks has resulted in hundreds of redundancies.}
{Đây là một danh từ dùng để chỉ sự cắt giảm biên chế hoặc tình trạng dư thừa nhân viên.}

\vocabentry{Resignation}
{/ˌrezɪɡˈneɪʃn/ (n)}
{An act of retiring or giving up a position.}
{The CEO announced his resignation following the financial scandal.}
{Đây là một danh từ dùng để chỉ sự từ chức hoặc đơn xin nghỉ việc.}

\vocabentry{Remuneration}
{/rɪˌmjuːnəˈreɪʃn/ (n)}
{Money paid for work or a service.}
{They offered a high level of remuneration to attract experienced candidates.}
{Đây là một danh từ trang trọng dùng để chỉ tiền thù lao, bao gồm lương và các khoản thưởng.}

\vocabentry{Promotion}
{/prəˈmoʊʃn/ (n)}
{The action of raising someone to a higher position or rank, or the publicizing of a product.}
{Hard work and dedication are the keys to earning a promotion.}
{Đây là một danh từ dùng để chỉ sự thăng tiến trong công việc hoặc hoạt động quảng bá}

\vocabentry{Apprentice}
{/əˈprentɪs/ (n)}
{A person who is learning a trade from a skilled employer.}
{He started his career as an apprentice in a local carpentry workshop.}
{Đây là một danh từ dùng để chỉ người học việc hoặc thực tập sinh nghề.}

\vocabentry{Commission}
{/kəˈmɪʃn/ (n)}
{An amount of money, typically a set percentage of the value involved, paid to an agent in a commercial transaction.}
{Sales agents receive a base salary plus a 5\% commission on every deal they close.}
{Đây là một danh từ dùng để chỉ tiền hoa hồng hưởng theo doanh số.}

\vocabentry{Telecommute}
{/ˌtelikəˈmjuːt/ (v)}
{To work from home, making use of the internet, email, and telephone.}
{More companies are allowing their employees to telecommute to improve work-life balance.}
{Đây là một động từ dùng để chỉ việc làm việc từ xa mà không cần đến văn phòng.}

\vocabentry{Incentive}
{/ɪnˈsentɪv/ (n)}
{A thing that motivates or encourages someone to do something, especially a payment or concession.}
{Cash bonuses are a common incentive for high-performing employees.}
{Đây là một danh từ dùng để chỉ động lực, phương tiện khích lệ nhân viên}

\vocabentry{Superannuation}
{/ˌsuːpərˌænjuˈeɪʃn/ (n)}
{Regular payment made into a fund by an employee towards a future pension.}
{A percentage of your monthly salary is automatically deducted for superannuation.}
{Đây là một danh từ trang trọng dùng để chỉ tiền đóng bảo hiểm hưu trí.}

\vocabentry{Dismissal}
{/dɪsˈmɪsl/ (n)}
{The act of ordering someone to leave a job.}
{Unfair dismissal is a serious issue that can lead to legal action against the employer.}
{Đây là một danh từ dùng để chỉ sự sa thải hoặc cho nghỉ việc.}

\vocabentry{Vocation}
{/voʊˈkeɪʃn/ (n)}
{A strong feeling of suitability for a particular career or occupation; a calling.}
{Many nurses feel that their work is a vocation rather than just a way to earn money.}
{Đây là một danh từ dùng để chỉ sự tập trung chuyên sâu vào một nghề nghiệp cụ thể, đam mê nghề nghiệp}

\vocabentry{Overtime}
{/ˈoʊvərtaɪm/ (n)}
{Time worked in addition to one's normal working hours.}
{Employees are paid double for any overtime worked during the weekends.}
{Đây là một danh từ dùng để chỉ thời gian làm thêm giờ.}

\vocabentry{White-collar}
{/ˌwaɪt ˈkɑːlər/ (adj)}
{Relating to the work done by people who work in offices.}
{The city has seen a significant increase in white-collar employment over the last decade.}
{Đây là một tính từ dùng để chỉ tầng lớp lao động trí óc hoặc nhân viên văn phòng.}

\vocabentry{Blue-collar}
{/ˌbluː ˈkɑːlər/ (adj)}
{Relating to manual work or workers.}
{Many blue-collar jobs in the manufacturing sector have been replaced by automation.}
{Đây là một tính từ dùng để chỉ tầng lớp lao động chân tay hoặc công nhân sản xuất.}

\vocabentry{Freelancer}
{/ˈfriːlænsər/ (n)}
{A person who works as a writer, designer, performer, or the like, selling work or services by the hour, day, job, etc.}
{As a freelancer, she has the flexibility to choose which projects she wants to work on.}
{Đây là một danh từ dùng để chỉ người làm nghề tự do.}

\vocabentry{Headhunter}
{/ˈhedhʌntər/ (n)}
{A person whose profession is to find executives or other highly skilled personnel and persuade them to change jobs.}
{A headhunter contacted him with a lucrative job offer from a rival company.}
{Đây là một danh từ dùng để chỉ người chuyên săn đầu người/tuyển dụng nhân sự cấp cao.}

\vocabentry{Workaholic}
{/ˌwɜːrkəˈhɑːlɪk/ (n)}
{A person who compulsively works excessively hard and long hours.}
{Being a workaholic often leads to burnout and strained personal relationships.}
{Đây là một danh từ dùng để chỉ người nghiện công việc.}

\vocabentry{Probation}
{/proʊˈbeɪʃn/ (n)}
{A period of time at the start of a new job when you are watched to see if you are suitable.}
{After a three-month probation period, he was offered a permanent contract.}
{Đây là một danh từ dùng để chỉ thời gian thử việc.}

\vocabentry{Bureaucracy}
{/bjʊˈrɑːkrəsi/ (n)}
{A system of government or management that is complicated and has too many rules.}
{The excessive bureaucracy in the company makes it difficult to get anything done quickly.}
{Đây là một danh từ dùng để chỉ bộ máy hành chính rườm rà.}

\vocabentry{Productivity}
{/ˌproʊdʌkˈtɪvəti/ (n)}
{The effectiveness of productive effort, especially in industry.}
{Investing in new technology can significantly boost the company's productivity.}
{Đây là một danh từ dùng để chỉ năng suất lao động.}

\vocabentry{Curriculum Vitae}
{/kəˌrɪkjələm ˈviːtaɪ/ (n)}
{A brief account of a person's education, qualifications, and previous experience.}
{You should tailor your curriculum vitae to highlight the skills relevant to the job.}
{Đây là một danh từ dùng để chỉ sơ yếu lý lịch/CV.}

\vocabentry{Fringe benefit}
{/ˌfrɪndʒ ˈbenɪfɪt/ (n)}
{An extra benefit supplementing an employee's salary.}
{Free health insurance and a company car are some of the fringe benefits offered by the firm.}
{Đây là một cụm danh từ dùng để chỉ các phúc lợi phụ thêm ngoài lương như bảo hiểm, xe cộ.}

\vocabentry{Labor-intensive}
{/ˈleɪbər ɪnˈtensɪv/ (adj)}
{(Of a process or industry) requiring a large amount of labor to produce its goods or services.}
{Agriculture in many developing countries remains highly labor-intensive.}
{Đây là một tính từ dùng để chỉ các ngành nghề cần nhiều sức lao động.}

\vocabentry{Seniority}
{/ˌsiːniˈɔːrəti/ (n)}
{The fact or state of being older or higher in position or status than someone else.}
{Promotions are often based on merit rather than just seniority.}
{Đây là một danh từ dùng để chỉ thâm niên hoặc cấp bậc cao trong công việc.}

\vocabentry{Severance pay}
{/ˈsevərəns peɪ/ (n)}
{Money paid by an employer to an employee whose job the employer has had to bring to an end.}
{He received six months' salary as severance pay after being laid off.}
{Đây là một cụm danh từ dùng để chỉ tiền trợ cấp thôi việc.}

\vocabentry{Workplace}
{/ˈwɜːrkpleɪs/ (n)}
{A place where people work, such as an office or factory.}
{Discrimination has no place in a modern, inclusive workplace.}
{Đây là một danh từ dùng để chỉ nơi làm việc.}

\vocabentry{Entrepreneur}
{/ˌɑːntrəprəˈnɜːr/ (n)}
{A person who organizes and operates a business or businesses, taking on greater than normal financial risks.}
{The young entrepreneur started his tech company in a small garage.}
{Đây là một danh từ dùng để chỉ doanh nhân/người khởi nghiệp.}

\vocabentry{Downsize}
{/ˈdaʊnˌsaɪz/ (v)}
{Make (a company or organization) smaller by eliminating staff positions.}
{To remain profitable, the airline had to downsize its workforce by 20\%.}
{Đây là một động từ dùng để chỉ việc cắt giảm quy mô nhân sự của một công ty.}

\vocabentry{Outsource}
{/ˈaʊtsɔːrs/ (v)}
{Obtain (goods or a service) from an outside or foreign supplier, especially in place of an internal source.}
{To reduce costs, the company decided to outsource its IT support.}
{Đây là một động từ dùng để chỉ việc thuê ngoài, chuyển công việc cho đơn vị bên ngoài}

\vocabentry{Personnel}
{/ˌpɜːrsəˈnel/ (n)}
{People employed in an organization or engaged in an organized undertaking.}
{All personnel are required to attend the safety training session tomorrow.}
{Đây là một danh từ dùng để chỉ nhân sự hoặc bộ phận nhân sự.}

\vocabentry{Remit}
{/ˈriːmɪt/ (n)}
{The task or area of activity officially assigned to an individual or organization.}
{Marketing falls within her remit, while financial planning is handled by the CFO.}
{Đây là một danh từ dùng để chỉ phạm vi trách nhiệm hoặc nhiệm vụ được giao.}

\vocabentry{Soft skills}
{/ˈsɔːft skɪlz/ (n)}
{Personal attributes that enable someone to interact effectively and harmoniously with other people.}
{In today's job market, soft skills are just as important as technical qualifications.}
{Đây là một cụm danh từ dùng để chỉ các kỹ năng mềm như giao tiếp, làm việc nhóm.}

\vocabentry{Appraisal}
{/əˈpreɪzl/ (n)}
{An act of assessing something or someone. (Đây là một danh từ dùng để chỉ sự đánh giá (thường là đánh giá hiệu suất nhân viên hàng năm).).}
{During the annual appraisal, employees discuss their goals and achievements with their managers.}
{}

\vocabentry{Dead-end}
{/ˌded ˈend/ (adj)}
{(Of a job) providing no prospects for promotion or advancement.}
{He felt stuck in a dead-end job and decided it was time for a career change.}
{Đây là một tính từ dùng để chỉ một công việc không có tương lai hoặc cơ hội thăng tiến.}

\vocabentry{Glass ceiling}
{/ˌɡlæs ˈsiːlɪŋ/ (n)}
{An unacknowledged barrier to advancement in a profession, especially affecting women and members of minorities. (Đây là một cụm danh từ dùng để chỉ rào cản vô hình ngăn cản sự thăng tiến (thường đối với phụ nữ).).}
{Many women in corporate roles still struggle to break through the glass ceiling.}
{}

\vocabentry{Job security}
{/ˈdʒɑːb sɪˈkjʊrəti/ (n)}
{The probability that an individual will keep their job.}
{In a volatile economy, job security has become a primary concern for many workers.}
{Đây là một cụm danh từ dùng để chỉ sự ổn định/an toàn của công việc.}

\vocabentry{Picket line}
{/ˈpɪkɪt laɪn/ (n)}
{A boundary established by workers on strike.}
{Union members stood on the picket line to protest against the proposed wage cuts.}
{Đây là một cụm danh từ dùng để chỉ hàng rào bãi công của công nhân.}

\vocabentry{Shortlist}
{/ˈʃɔːrtlɪst/ (v)}
{Put (someone or something) on a shortlist.}
{From over a hundred applicants, only five were shortlisted for the final interview.}
{Đây là một động từ dùng để chỉ việc đưa ai đó vào danh sách sơ tuyển/vòng rút gọn.}

\vocabentry{Unemployment}
{/ˌʌnɪmˈplɔɪmənt/ (n)}
{The state of being unemployed.}
{High levels of youth unemployment can lead to social instability.}
{Đây là một danh từ dùng để chỉ tình trạng thất nghiệp.}

\vocabentry{Wage}
{/weɪdʒ/ (n)}
{A fixed regular payment, typically paid on a daily or weekly basis, made by an employer to an employee.}
{The company offers a competitive wage and a range of fringe benefits.}
{Đây là một danh từ dùng để chỉ lương cố định, thường trả theo ngày hoặc theo tuần}

\vocabentry{Work-life balance}
{/ˌwɜːrk laɪf ˈbæləns/ (n)}
{The division of one's time and focus between working and family or leisure activities.}
{Flexible working hours can significantly improve an employee's work-life balance.}
{Đây là một cụm danh từ dùng để chỉ sự cân bằng giữa công việc và đời sống riêng.}

\vocabentry{Full-time}
{/ˌfʊl ˈtaɪm/ (adj)}
{Occupying the whole of someone's available working time.}
{She recently transitioned from a part-time role to a full-time position.}
{Đây là một tính từ dùng để chỉ làm việc toàn thời gian.}

\vocabentry{Part-time}
{/ˌpɑːrt ˈtaɪm/ (adj)}
{For only part of the usual working day or week.}
{Many students take part-time jobs to help pay for their university fees.}
{Đây là một tính từ dùng để chỉ làm việc bán thời gian.}

\vocabentry{Maternity leave}
{/məˈtɜːrnəti liːv/ (n)}
{A period of absence from work allowed to a mother before and after the birth of her child.}
{She will be on maternity leave for six months starting next January.}
{Đây là một cụm danh từ dùng để chỉ kỳ nghỉ thai sản của người mẹ.}

\vocabentry{Paternity leave}
{/pəˈtɜːrnəti liːv/ (n)}
{A period of absence from work allowed to a father after the birth of his child.}
{The new policy allows fathers to take up to two weeks of paid paternity leave.}
{Đây là một cụm danh từ dùng để chỉ kỳ nghỉ thai sản dành cho người cha.}

\vocabentry{Labor union}
{/ˈleɪbər ˈjuːniən/ (n)}
{An organized association of workers formed to protect and further their rights and interests.}
{The labor union is negotiating with the management for better working conditions.}
{Đây là một cụm danh từ dùng để chỉ công đoàn lao động.}

\vocabentry{Flexitime}
{/ˈfleksitaɪm/ (n)}
{A system of working a set number of hours with the starting and finishing times chosen by the employee.}
{Flexitime allows parents to drop their children at school before heading to the office.}
{Đây là một danh từ dùng để chỉ hệ thống giờ làm việc linh hoạt.}

\vocabentry{Job satisfaction}
{/ˈdʒɑːb ˌsætɪsˈfækʃn/ (n)}
{A feeling of fulfillment or enjoyment that a person derives from their job.}
{Money is important, but job satisfaction is the key to long-term happiness at work.}
{Đây là một cụm danh từ dùng để chỉ sự hài lòng trong công việc.}

\vocabentry{Stonewall}
{/ˈstoʊnwɔːl/ (v)}
{To delay or block (a process or person) by refusing to answer questions or by giving evasive replies.}
{Management was accused of stonewalling during the discussions about salary increases.}
{Đây là một động từ dùng để chỉ việc trì hoãn hoặc cản trở quá trình đàm phán trong công việc.}

\vocabentry{Contractor}
{/ˈkɑːntræktər/ (n)}
{A person or company that undertakes a contract to provide materials or labor to perform a service or do a job.}
{The government hired an independent contractor to oversee the bridge construction.}
{Đây là một danh từ dùng để chỉ nhà thầu hoặc người làm việc theo hợp đồng dự án.}

\vocabentry{Mentoring}
{/ˈmentɔːrɪŋ/ (n)}
{The practice of helping and advising a less experienced person.}
{Mentoring programs are a great way to pass on knowledge to junior staff members.}
{Đây là một danh từ dùng để chỉ hoạt động cố vấn, hướng dẫn người mới.}

\vocabentry{Overworked}
{/ˌoʊvərˈwɜːrkt/ (adj)}
{(Of a person) exhausted by too much work.}
{Overworked nurses often struggle to maintain the highest standards of patient care.}
{Đây là một tính từ dùng để chỉ tình trạng làm việc quá sức.}

\vocabentry{Sick leave}
{/ˈsɪk liːv/ (n)}
{Leave of absence from work through illness.}
{You need to provide a doctor's certificate if you are on sick leave for more than three days.}
{Đây là một cụm danh từ dùng để chỉ kỳ nghỉ ốm.}

\vocabentry{Turnover}
{/ˈtɜːrnoʊvər/ (n)}
{The rate at which employees leave a workforce and are replaced.}
{High employee turnover is often a sign of poor management or low morale.}
{Đây là một danh từ dùng để chỉ tỷ lệ luân chuyển/thay thế nhân viên.}

\vocabentry{Blue-chip}
{/ˌbluː ˈtʃɪp/ (adj)}
{Denoting companies or their shares considered to be a reliable investment.}
{Getting a job at a blue-chip company like Google is highly competitive.}
{Đây là một tính từ dùng để chỉ những công ty lớn, uy tín và ổn định về tài chính.}

\vocabentry{Clerical}
{/ˈklerɪkl/ (adj)}
{Relating to office work.}
{She started as a clerical assistant and eventually became the office manager.}
{Đây là một tính từ dùng để chỉ các công việc văn phòng liên quan đến giấy tờ, sổ sách.}

\vocabentry{Empowerment}
{/ɪmˈpaʊərmənt/ (n)}
{The process of becoming stronger and more confident, especially in controlling one's life and claiming one's rights.}
{Employee empowerment leads to more innovation and better problem-solving within teams.}
{Đây là một danh từ dùng để chỉ sự trao quyền cho nhân viên.}

\vocabentry{Grievance}
{/ˈɡriːvəns/ (n)}
{A real or imagined wrong or other cause for complaint or protest, especially unfair treatment.}
{If you have a grievance, you should follow the formal procedure outlined in the employee handbook.}
{Đây là một danh từ dùng để chỉ sự khiếu nại hoặc nỗi bất bình của nhân viên.}

\vocabentry{Hierarchy}
{/ˈhaɪərɑːrki/ (n)}
{A system or organization in which people or groups are ranked one above the other according to status or authority.}
{The traditional corporate hierarchy is being replaced by more flexible, flat structures.}
{Đây là một danh từ dùng để chỉ hệ thống cấp bậc trong tổ chức.}

\vocabentry{Internship}
{/ˈɪntɜːrnʃɪp/ (n)}
{The position of a student or trainee who works in an organization, sometimes without pay, in order to gain work experience.}
{Doing an internship is a great way to gain practical experience and network in your field.}
{Đây là một danh từ dùng để chỉ kỳ thực tập.}

\vocabentry{Knowledge-based}
{/ˈnɑːlɪdʒ beɪst/ (adj)}
{(Of an economy or industry) driven by the use of knowledge.}
{Most developed nations are transitioning towards a knowledge-based economy.}
{Đây là một tính từ dùng để chỉ các ngành kinh tế dựa trên tri thức.}

\vocabentry{Lucrative}
{/ˈluːkrətɪv/ (adj)}
{Producing a great deal of profit.}
{He left his job in the public sector for a more lucrative position in finance.}
{Đây là một tính từ dùng để chỉ một công việc hoặc dự án mang lại lợi nhuận cao, hái ra tiền.}

\vocabentry{Manual labor}
{/ˈmænjuəl ˈleɪbər/ (n)}
{Physical work done by people.}
{Despite the rise of robots, some tasks still require skilled manual labor.}
{Đây là một cụm danh từ dùng để chỉ lao động chân tay.}

\vocabentry{Moonlight}
{/ˈmuːnlaɪt/ (v)}
{Have a second job in addition to one's main employment, typically without the knowledge of one's main employer.}
{He is moonlighting as a taxi driver to pay off his debts.}
{Đây là một động từ dùng để chỉ việc làm thêm nghề tay trái vào ban đêm hoặc ngoài giờ chính.}

\vocabentry{Networking}
{/ˈnetwɜːrkɪŋ/ (n)}
{The action or process of interacting with others to exchange information and develop professional or social contacts.}
{Networking is often the most effective way to find new job opportunities.}
{Đây là một danh từ dùng để chỉ việc thiết lập mạng lưới quan hệ chuyên môn.}

\vocabentry{Occupational hazard}
{/ˌɑːkjuˈpeɪʃənl ˈhæzərd/ (n)}
{A risk or danger that is part of a particular job.}
{Back pain is a common occupational hazard for people who sit at a desk all day.}
{Đây là một cụm danh từ dùng để chỉ những mối nguy hiểm đặc thù trong nghề nghiệp.}

\vocabentry{Perk}
{/pɜːrk/ (n)}
{A benefit to which one is entitled because of one's job. (Đây là một danh từ (không trang trọng) dùng để chỉ các đặc quyền, bổng lộc phụ thêm.).}
{One of the perks of working for an airline is free or discounted flights.}
{}

\vocabentry{Qualifications}
{/ˌkwɑːlɪfɪˈkeɪʃnz/ (n)}
{A pass of an examination or an official completion of a course, especially one conferring status as a recognized practitioner.}
{Do you have any professional qualifications in accounting or finance?}
{Đây là một danh từ dùng để chỉ trình độ chuyên môn, bằng cấp nghề nghiệp}

\vocabentry{Recruiter}
{/rɪˈkruːtər/ (n)}
{A person whose job is to find people to work for a company.}
{The recruiter was impressed by his experience and strong communication skills.}
{Đây là một danh từ dùng để chỉ nhà tuyển dụng.}

\vocabentry{Self-employed}
{/ˌself ɪmˈplɔɪd/ (adj)}
{Working for oneself as a freelancer or the owner of a business rather than for an employer.}
{Being self-employed offers freedom but also comes with financial uncertainty.}
{Đây là một tính từ dùng để chỉ tình trạng tự làm chủ.}

\vocabentry{Skill set}
{/ˈskɪl set/ (n)}
{A person's range of skills or abilities.}
{She has a diverse skill set that includes coding, graphic design, and project management.}
{Đây là một cụm danh từ dùng để chỉ bộ kỹ năng hoặc phạm vi năng lực của một người.}

\vocabentry{Tax return}
{/ˈtæks rɪˈtɜːrn/ (n)}
{A form on which a taxpayer makes an annual statement of income and personal circumstances.}
{You must submit your tax return by the end of the financial year to avoid a penalty.}
{Đây là một cụm danh từ dùng để chỉ việc kê khai thuế thu nhập.}

\vocabentry{Underemployed}
{/ˌʌndərɪmˈplɔɪd/ (adj)}
{(Of a person) not having enough paid work or not doing work that makes full use of their skills and abilities.}
{Many PhD graduates are currently underemployed, working in roles that do not require their high level of expertise.}
{Đây là một tính từ dùng để chỉ tình trạng làm việc dưới năng lực hoặc không đủ thời gian.}

\vocabentry{Vacancy}
{/ˈveɪkənsi/ (n)}
{An unoccupied position or job that is available to be filled.}
{We have several vacancies for experienced software engineers.}
{Đây là một danh từ dùng để chỉ vị trí tuyển dụng, công việc đang bỏ trống}

\vocabentry{Whistleblower}
{/ˈwɪslbloʊər/ (n)}
{A person who informs on a person or organization engaged in an illicit activity.}
{The whistleblower exposed the company's illegal waste disposal practices.}
{Đây là một danh từ dùng để chỉ người tố giác hành vi sai trái trong nội bộ tổ chức.}

\vocabentry{Xenophobia}
{/ˌzenəˈfoʊbiə/ (n)}
{Dislike of or prejudice against people from other countries.}
{HR policies must be strictly enforced to prevent any form of xenophobia in the workplace.}
{Đây là một danh từ có thể liên quan đến vấn đề phân biệt đối xử với lao động nước ngoài tại nơi làm việc.}

\vocabentry{Yield}
{/jiːld/ (v)}
{Produce or provide (a natural, agricultural, or industrial product); generate (a financial return).}
{Investing in new technology will yield significant benefits for the company.}
{Đây là một động từ dùng để chỉ sự mang lại kết quả, tạo ra lợi nhuận hoặc sản phẩm}

\vocabentry{Zero-hours contract}
{/ˌzɪroʊ aʊərz ˈkɑːntrækt/ (n)}
{A type of contract between an employer and a worker where the employer is not obliged to provide any minimum working hours.}
{Zero-hours contracts offer flexibility for employers but provide little security for workers.}
{Đây là một cụm danh từ dùng để chỉ hợp đồng không cam kết giờ làm việc tối thiểu.}

\vocabentry{Accredit}
{/əˈkredɪt/ (v)}
{Give official authorization to or approval of.}
{The institution is accredited by the national board of education.}
{Đây là một động từ dùng để chỉ việc cấp phép hoặc công nhận chính thức năng lực nghề nghiệp.}

\vocabentry{Board of directors}
{/ˌbɔːrd əv dəˈrektərz/ (n)}
{A group of people who jointly supervise the activities of an organization.}
{The board of directors met yesterday to discuss the company's new expansion strategy.}
{Đây là một cụm danh từ dùng để chỉ hội đồng quản trị.}

\vocabentry{Candidate}
{/ˈkændɪdət/ (n)}
{A person who applies for a job or is nominated for election.}
{We interviewed several excellent candidates, but none of them quite met our requirements.}
{Đây là một danh từ dùng để chỉ ứng viên.}

\vocabentry{Deadline}
{/ˈdedlaɪn/ (n)}
{The latest time or date by which something should be completed.}
{We are working hard to meet the deadline for the final report.}
{Đây là một danh từ dùng để chỉ hạn chót công việc.}

\vocabentry{Escalate}
{/ˈeskəleɪt/ (v)}
{Increase rapidly; refer a situation to a higher authority.}
{If the customer is still not satisfied, you should escalate the issue to your supervisor.}
{Đây là một động từ dùng để chỉ việc báo cáo vấn đề lên cấp trên để giải quyết.}

\vocabentry{Flowchart}
{/ˈfloʊtʃɑːrt/ (n)}
{A diagram of the sequence of movements or actions which is part of a complex system or activity.}
{The manager used a flowchart to explain the new production process to the team.}
{Đây là một danh từ dùng để chỉ sơ đồ quy trình công việc.}

\vocabentry{Gap year}
{/ˈɡæp jɪər/ (n)}
{A period, typically an academic year, taken by a student as a break between secondary school and higher education. (Đây là một cụm danh từ dùng để chỉ năm nghỉ giữa các bậc học (thường để đi làm lấy kinh nghiệm).).}
{He spent his gap year working as a volunteer in Africa.}
{}

\vocabentry{Handover}
{/ˈhændoʊvər/ (n)}
{The act of passing on authority or responsibility to another person.}
{There will be a two-week handover period before the old manager leaves.}
{Đây là một danh từ dùng để chỉ sự bàn giao công việc.}

\vocabentry{Intern}
{/ˈɪntɜːrn/ (n)}
{A student or trainee who works, sometimes without pay, in order to gain work experience.}
{The company hires several interns every summer to assist with administrative tasks.}
{Đây là một danh từ dùng để chỉ thực tập sinh.}

\vocabentry{Job description}
{/ˈdʒɑːb dɪˌskrɪpʃn/ (n)}
{A formal account of an employee's responsibilities.}
{Please read the job description carefully before applying for the position.}
{Đây là một cụm danh từ dùng để chỉ bản mô tả công việc.}

\vocabentry{Key performance indicator}
{/ˌkiː pərˈfɔːrməns ˌɪndɪkeɪtər/ (n)}
{A quantifiable measure used to evaluate the success of an organization or of a particular activity in which it engages.}
{Sales growth is a key performance indicator for our marketing team.}
{Đây là một cụm danh từ dùng để chỉ chỉ số đánh giá hiệu quả công việc - KPI.}

\vocabentry{Liaison}
{/liˈeɪzɑːn/ (n)}
{Communication or cooperation which facilitates a close working relationship between people or organizations.}
{She acts as a liaison between the engineering department and the sales team.}
{Đây là một danh từ dùng để chỉ sự liên lạc hoặc người giữ vai trò kết nối giữa các bên.}

\vocabentry{Managerial}
{/ˌmænəˈdʒɪriəl/ (adj)}
{Relating to management or managers.}
{He has several years of managerial experience in the retail industry.}
{Đây là một tính từ dùng để chỉ những gì thuộc về quản lý.}

\vocabentry{Notice period}
{/ˈnoʊtɪs ˌpɪriəd/ (n)}
{The time between the date an employee gives notice and the date they leave their job.}
{My contract states that I must give a one-month notice period before resigning.}
{Đây là một cụm danh từ dùng để chỉ thời gian thông báo nghỉ việc trước khi chính thức rời đi.}

\vocabentry{Objective}
{/əbˈdʒektɪv/ (n)}
{A thing aimed at or sought; a goal.}
{Our primary objective this year is to increase our market share by 10\%.}
{Đây là một danh từ dùng để chỉ mục tiêu công việc.}

\vocabentry{Portfolio}
{/pɔːrtˈfoʊlioʊ/ (n)}
{A set of pieces of creative work collected by someone to display their skills.}
{Graphic designers should have a strong portfolio to show to potential clients.}
{Đây là một danh từ dùng để chỉ hồ sơ năng lực/tập hợp các sản phẩm đã thực hiện.}

\vocabentry{Quota}
{/ˈkwoʊtə/ (n)}
{A fixed share of something that a person or group is entitled to receive or is bound to contribute.}
{The sales team is struggling to meet their monthly sales quota.}
{Đây là một danh từ dùng để chỉ chỉ tiêu hoặc hạn ngạch công việc.}

\vocabentry{Retund}
{/rɪˈtɜːrn/ (v)}
{Come or go back to a place or person.}
{She plans to return to work after her maternity leave finishes.}
{Trong công việc, có thể dùng chỉ việc quay lại làm việc - return to work.}

\vocabentry{Stakeholder}
{/ˈsteɪkhoʊldər/ (n)}
{A person with an interest or concern in something, especially a business.}
{We need to consider the interests of all stakeholders before making a final decision.}
{Đây là một danh từ dùng để chỉ các bên liên quan trong một dự án hoặc công ty.}

\vocabentry{Task-oriented}
{/ˈtæsk ˈɔːriəntɪd/ (adj)}
{Focusing on the completion of specific tasks.}
{A task-oriented approach can be very effective for meeting tight deadlines.}
{Đây là một tính từ dùng để chỉ phong cách làm việc tập trung vào nhiệm vụ cụ thể.}

\vocabentry{Undertaking}
{/ˌʌndərˈteɪkɪŋ/ (n)}
{A formal pledge or promise to do something; a task.}
{Building the new highway is a massive undertaking that will take several years.}
{Đây là một danh từ dùng để chỉ một công việc hoặc cam kết thực hiện một nhiệm vụ.}

\vocabentry{Verification}
{/ˌverɪfɪˈkeɪʃn/ (n)}
{The process of establishing the truth, accuracy, or validity of something. (Đây là một danh từ dùng để chỉ sự xác minh (như xác minh bằng cấp).).}
{The HR department is currently conducting a background verification of the new hires.}
{}

\vocabentry{Workforce}
{/ˈwɜːrkfɔːrs/ (n)}
{The people engaged in or available for work, either in a country or area or in a particular company.}
{A diverse workforce can bring a wider range of ideas and perspectives to a company.}
{Đây là một danh từ dùng để chỉ lực lượng lao động, nguồn nhân lực}

\vocabentry{Zero tolerance}
{/ˌzɪroʊ ˈtɑːlərəns/ (n)}
{Refusal to accept antisocial behavior, typically by strict and uncompromising application of the law. (Đây là một cụm danh từ dùng để chỉ chính sách không khoan nhượng (ví dụ đối với quấy rối nơi làm việc).).}
{The company has a zero tolerance policy towards bullying and harassment.}
{}

\vocabentry{Absence}
{/ˈæbsəns/ (n)}
{The state of being away from a place or person.}
{He has had several unexcused absences this month.}
{Đây là một danh từ dùng để chỉ sự vắng mặt tại nơi làm việc.}

\vocabentry{Benchmark}
{/ˈbentʃmɑːrk/ (n)}
{A standard or point of reference against which things may be compared or assessed.}
{Our sales performance is often used as a benchmark for the rest of the industry.}
{Đây là một danh từ dùng để chỉ điểm chuẩn hoặc tiêu chuẩn đối đối chiếu hiệu quả công việc.}

\vocabentry{Collaborate}
{/kəˈlæbəreɪt/ (v)}
{Work jointly on an activity, especially to produce or create something.}
{Designers and engineers need to collaborate closely to create a successful product.}
{Đây là một động từ dùng để chỉ việc cộng tác/phối hợp làm việc.}

\vocabentry{Diversify}
{/daɪˈvɜːrsɪfaɪ/ (v)}
{Make or become more diverse or varied.}
{The company is trying to diversify its business to reduce risk.}
{Đây là một động từ dùng để chỉ việc đa dạng hóa kỹ năng hoặc lĩnh vực hoạt động.}

\vocabentry{Escalation}
{/ˌeskəˈleɪʃn/ (n)}
{A rapid increase in size, scope, or intensity; the process of referring a matter to a higher authority.}
{The escalation of the conflict led to a total breakdown in communications.}
{Đây là một danh từ dùng để chỉ sự báo cáo vượt cấp khi có sự cố.}

\vocabentry{Facilitate}
{/fəˈsɪlɪteɪt/ (v)}
{Make (an action or process) easy or easier.}
{The new software is designed to facilitate communication between different departments.}
{Đây là một động từ dùng để chỉ việc tạo điều kiện thuận lợi cho công việc.}

\vocabentry{Goal-oriented}
{/ˈɡoʊl ˈɔːriəntɪd/ (adj)}
{Focusing on achieving specific goals.}
{We are looking for a goal-oriented individual who can drive sales growth.}
{Đây là một tính từ dùng để chỉ phong cách làm việc tập trung vào mục tiêu.}

\vocabentry{Human resources}
{/ˌhjuːmən ˈriːsɔːrsɪz/ (n)}
{The department of a business or organization that deals with the hiring, administration, and training of personnel.}
{Please contact the human resources department if you have any questions about your benefits.}
{Đây là một cụm danh từ dùng để chỉ bộ phận nhân sự - HR.}

\vocabentry{Incentivize}
{/ɪnˈsentɪvaɪz/ (v)}
{Provide (someone) with an incentive for doing something.}
{The government is trying to incentivize businesses to invest in renewable energy.}
{Đây là một động từ dùng để chỉ việc khuyến khích/động viên bằng các phần thưởng.}

\vocabentry{Junior}
{/ˈdʒuːniər/ (adj)}
{Denoting low or lower status or rank in an organization.}
{He started his career as a junior developer and was promoted within a year.}
{Đây là một tính từ dùng để chỉ cấp bậc nhân viên mới hoặc cấp thấp.}

\vocabentry{Knowledgeable}
{/ˈnɑːlɪdʒəbl/ (adj)}
{Intelligent and well informed; having or showing knowledge or intelligence.}
{Our customer service staff are very knowledgeable about all of our services.}
{Đây là một tính từ dùng để chỉ người có kiến thức sâu rộng, hiểu biết}

\vocabentry{Livelihood}
{/ˈlaɪvlihʊd/ (n)}
{A means of securing the necessities of life.}
{Thousands of people lost their livelihood when the factory closed down.}
{Đây là một danh từ dùng để chỉ kế sinh nhai/nguồn sống.}

\vocabentry{Mediation}
{/ˌmiːdiˈeɪʃn/ (n)}
{Intervention in a dispute in order to resolve it.}
{The dispute was finally settled through a process of independent mediation.}
{Đây là một danh từ dùng để chỉ sự hòa giải mâu thuẫn lao động.}

\vocabentry{Non-profit}
{/ˌnɑːn ˈprɑːfɪt/ (adj)}
{(Of an organization) not established for the purpose of making a profit.}
{She decided to leave the corporate world to work for a non-profit organization.}
{Đây là một tính từ dùng để chỉ các tổ chức phi lợi nhuận.}

\vocabentry{Outplacement}
{/ˈaʊtˌpleɪsmənt/ (n)}
{The provision of assistance to redundant employees in finding new employment.}
{The company offered an outplacement service to help laid-off staff update their CVs and find new roles.}
{Đây là một danh từ dùng để chỉ dịch vụ hỗ trợ tìm việc mới cho nhân viên bị cắt giảm biên chế.}

\vocabentry{Payroll}
{/ˈpeɪroʊl/ (n)}
{A list of a company's employees and the amount of money they are to be paid.}
{The company has over five hundred employees on its payroll.}
{Đây là một danh từ dùng để chỉ bảng lương hoặc danh sách trả lương.}

\vocabentry{Quarterly}
{/ˈkwɔːrtərli/ (adv)}
{At three-monthly intervals.}
{We meet quarterly to review our progress and set new goals.}
{Đây là một trạng từ thường dùng để chỉ các đợt đánh giá hoặc báo cáo hàng quý.}

\vocabentry{Redeploy}
{/ˌriːdɪˈplɔɪ/ (v)}
{Assign (troops, employees, or resources) to a new place or task.}
{When the project was canceled, the staff were redeployed to other departments.}
{Đây là một động từ dùng để chỉ việc điều động lại nhân sự sang vị trí hoặc bộ phận khác.}

\vocabentry{Supervisor}
{/ˈsuːpərvaɪzər/ (n)}
{A person who supervises a person or an activity, especially workers.}
{My supervisor gave me some very helpful feedback on my last project.}
{Đây là một danh từ dùng để chỉ người giám sát, quản lý trực tiếp}

\vocabentry{Teamwork}
{/ˈtiːmwɜːrk/ (n)}
{The combined action of a group of people, especially when effective and efficient.}
{Success in this role requires excellent communication and teamwork skills.}
{Đây là một danh từ dùng để chỉ sự làm việc nhóm.}

\vocabentry{Understaffed}
{/ˌʌndərˈstæft/ (adj)}
{Having fewer staff than is needed.}
{The hospital is severely understaffed, which is putting a lot of pressure on the existing nurses.}
{Đây là một tính từ dùng để chỉ tình trạng thiếu hụt nhân sự.}

\vocabentry{Voluntary}
{/ˈvɑːlənteri/ (adj)}
{Done, given, or acting of one's own free will.}
{Several employees decided to take voluntary redundancy.}
{Đây là một tính từ có thể dùng chỉ việc tình nguyện nghỉ việc hoặc làm việc không lương.}

\vocabentry{Workshop}
{/ˈwɜːrkʃɑːp/ (n)}
{A meeting at which a group of people engage in intensive discussion and activity on a particular subject or project.}
{The HR department is organizing a workshop on effective time management.}
{Đây là một danh từ dùng để chỉ buổi hội thảo thực hành hoặc tập huấn.}

\vocabentry{Zone}
{/zoʊn/ (n)}
{An area or stretch of land having a particular characteristic, purpose, or use.}
{The new factory is located in a special economic zone with lower tax rates.}
{Trong công việc có thể là khu công nghiệp - industrial zone.}

\vocabentry{Affiliate}
{/əˈfɪlieɪt/ (n)}
{A person or organization officially attached to a larger body.}
{He works for an affiliate of a large multinational corporation.}
{Đây là một danh từ dùng để chỉ chi nhánh hoặc đối tác liên kết.}

\vocabentry{Beneficiary}
{/ˌbenɪˈfɪʃieri/ (n)}
{A person who derives advantage from something, especially a trust, will, or life insurance policy.}
{Employees are the primary beneficiaries of the new health and wellness program.}
{Đây là một danh từ có thể chỉ người hưởng lợi từ các chính sách bảo hiểm công ty.}

\vocabentry{Clientele}
{/ˌklaɪənˈtel/ (n)}
{The customers of a shop, bar, or place of entertainment.}
{The firm has a high-end clientele that includes many major corporations.}
{Đây là một danh từ dùng để chỉ tệp khách hàng của một doanh nghiệp.}

\vocabentry{Delegation}
{/ˌdelɪˈɡeɪʃn/ (n)}
{The act of giving a particular job, duty, right, etc. to someone else so that they do it for you.}
{Effective delegation is a key skill for any successful manager.}
{Đây là một danh từ dùng để chỉ sự ủy quyền/giao phó trách nhiệm.}

\vocabentry{Executive}
{/ɪɡˈzekjətɪv/ (n)}
{A person with senior managerial responsibility in a business organization.}
{She is a senior executive at one of the world's largest advertising agencies.}
{Đây là một danh từ dùng để chỉ giám đốc điều hành hoặc nhân sự cấp cao.}

\vocabentry{Fulfillment}
{/fʊlˈfɪlmənt/ (n)}
{Satisfaction or happiness as a result of fully developing one's abilities or character.}
{He finds a lot of personal fulfillment in helping others achieve their goals.}
{Đây là một danh từ dùng để chỉ cảm giác thành tựu/mãn nguyện trong công việc.}

\vocabentry{Graduate}
{/ˈɡrædʒuət/ (n)}
{A person who has successfully completed a course of study or training, especially a person who has been awarded an undergraduate or first academic degree.}
{We are looking for bright and ambitious graduates to join our management training program.}
{Thường dùng trong tuyển dụng - graduate recruitment.}

\vocabentry{Headquarters}
{/ˈhedkwɔːrtərz/ (n)}
{The premises serving as the center of an organization, from which military or business operations are controlled.}
{The company's headquarters are located in New York City.}
{Đây là một danh từ dùng để chỉ trụ sở chính.}

\vocabentry{Industrious}
{/ɪnˈdʌstriəs/ (adj)}
{Diligent and hard-working.}
{She is an industrious student who always finishes her assignments on time.}
{Đây là một tính từ dùng để chỉ một người làm việc cần cù, siêng năng.}

\vocabentry{Job-share}
{/ˈdʒɑːb ʃer/ (v)}
{(Of two people) share the work and pay of a single full-time job.}
{They decided to job-share so they could both spend more time with their families.}
{Đây là một động từ dùng để chỉ việc chia sẻ một công việc toàn thời gian cho hai người làm.}

\vocabentry{Keynote}
{/ˈkiːnoʊt/ (n)}
{A central theme used in a speech or program.}
{The CEO gave an inspiring keynote address at the annual company conference.}
{Thường dùng chỉ bài phát biểu chính trong hội nghị - keynote speech.}

\vocabentry{Labor market}
{/ˈleɪbər ˈmɑːrkɪt/ (n)}
{The supply of people in a particular country or area who are able and willing to work in relation to the number of jobs available.}
{The labor market is becoming increasingly competitive for recent graduates.}
{Đây là một cụm danh từ dùng để chỉ thị trường lao động.}

\vocabentry{Misconduct}
{/ˌmɪsˈkɑːndʌkt/ (n)}
{Unacceptable or improper behavior, especially by a professional person.}
{He was dismissed for professional misconduct following an internal investigation.}
{Đây là một danh từ dùng để chỉ hành vi sai trái/vi phạm kỷ luật tại nơi làm việc.}

\vocabentry{Negotiate}
{/nəˈɡoʊʃieɪt/ (v)}
{Try to reach an agreement or compromise by discussion with others.}
{You should always try to negotiate your salary before accepting a new job offer.}
{Đây là một động từ dùng để chỉ việc thương lượng, đàm phán lương bổng hoặc hợp đồng.}

\vocabentry{Orientation}
{/ˌɔːriənˈteɪʃn/ (n)}
{A program of introduction for newcomers to a college or other institution.}
{The orientation day helped new employees get used to the office and meet their colleagues.}
{Đây là một danh từ dùng để chỉ buổi định hướng/hướng dẫn cho nhân viên mới.}

\vocabentry{Profitability}
{/ˌprɑːfɪtəˈbɪləti/ (n)}
{The degree to which a business or activity yields profit or financial gain.}
{The new CEO's main goal is to improve the company's profitability.}
{Đây là một danh từ dùng để chỉ khả năng sinh lời.}

\vocabentry{Qualification}
{/ˌkwɑːlɪfɪˈkeɪʃn/ (n)}
{A pass of an examination or an official record of achievement, especially in some specified field.}
{She has all the right qualifications and experience for the job.}
{Đây là một danh từ dùng để chỉ năng lực/bằng cấp.}

\vocabentry{Retirement}
{/rɪˈtaɪərmənt/ (n)}
{The action or fact of leaving one's job and ceasing to work.}
{He is looking forward to his retirement so he can travel more.}
{Đây là một danh từ dùng để chỉ sự nghỉ hưu.}

\vocabentry{Shift work}
{/ˈʃɪft wɜːrk/ (n)}
{Work scheduled in a system of shifts.}
{Shift work can be difficult because it disrupts your sleep pattern.}
{Đây là một cụm danh từ dùng để chỉ công việc làm theo ca.}

\vocabentry{Task}
{/tæsk/ (n)}
{A piece of work to be done or undertaken.}
{Each team member was assigned a specific task to complete by the end of the day.}
{Đây là một danh từ dùng để chỉ nhiệm vụ/công việc cụ thể.}

\vocabentry{Underprivileged}
{/ˌʌndərˈprɪvəlɪdʒd/ (adj)}
{Not enjoying the same standard of living or rights as the majority of people in a society.}
{The company has a program to provide training and jobs for underprivileged youth.}
{Có thể liên quan đến các chương trình hỗ trợ việc làm.}

\vocabentry{Wage earner}
{/ˈweɪdʒ ˌɜːrnər/ (n)}
{A person who works for wages.}
{In many families, both parents are now wage earners.}
{Đây là một cụm danh từ dùng để chỉ người làm công ăn lương.}

\vocabentry{Whistleblowing}
{/ˈwɪslbloʊɪŋ/ (n)}
{The act of informing on a person or organization engaged in an illicit activity.}
{Whistleblowing is essential for maintaining transparency and ethical standards in business.}
{Đây là một danh từ dùng để chỉ hành động tố giác sai phạm.}

\vocabentry{Yielding}
{/ˈjiːldɪŋ/ (adj)}
{Producing a result.}
{A more flexible approach to management is yielding positive results for the company.}
{Có thể dùng chỉ hiệu quả công việc.}

\vocabentry{Zeal}
{/ziːl/ (n)}
{Great energy or enthusiasm in pursuit of a cause or an objective.}
{She approached her new role with great zeal and dedication.}
{Đây là một danh từ dùng để chỉ lòng nhiệt huyết trong công việc.}

\vocabentry{Advancement}
{/ədˈvænsmənt/ (n)}
{The process of promoting a cause or plan.}
{There are many opportunities for career advancement within the company.}
{Thường dùng trong thăng tiến sự nghiệp - career advancement.}

\vocabentry{Background check}
{/ˈbækɡraʊnd tʃek/ (n)}
{A process a person or company uses to verify that a person is who they say they are.}
{All new employees must undergo a thorough background check before they start work.}
{Đây là một cụm danh từ dùng để chỉ việc kiểm tra lý lịch nhân viên.}

\vocabentry{Competency}
{/ˈkɑːmpɪtənsi/ (n)}
{The ability to do something successfully or efficiently.}
{Core competencies are the essential skills needed for a particular role.}
{Đây là một danh từ dùng để chỉ năng lực/khả năng thành thạo.}

\vocabentry{Discipline}
{/ˈdɪsəplɪn/ (n)}
{The practice of training people to obey rules or a code of behavior.}
{Maintaining a high level of discipline is essential for a safe and productive workplace.}
{Trong công việc là kỷ luật lao động.}

\vocabentry{Efficient}
{/ɪˈfɪʃnt/ (adj)}
{(Of a system or machine) achieving maximum productivity with minimum wasted effort or expense.}
{The new production process is much more efficient than the old one.}
{Đây là một tính từ dùng để chỉ tính hiệu quả.}

\vocabentry{Full-stack}
{/ˌfʊl ˈstæk/ (adj)}
{Relating to or being a computer programmer who is able to work with both the front end and back end of a website or application.}
{The company is looking for a full-stack developer to lead their new tech project.}
{Đây là một thuật ngữ chuyên ngành IT phổ biến trong tuyển dụng hiện nay, chỉ lập trình viên toàn diện}

\vocabentry{Guideline}
{/ˈɡaɪdlaɪn/ (n)}
{A general rule, principle, or piece of advice.}
{Please follow the safety guidelines carefully when using the machinery.}
{Trong công việc là các hướng dẫn thực hiện nhiệm vụ.}

\vocabentry{Hire}
{/ˈhaɪər/ (v)}
{Employ (someone) for wages.}
{We need to hire more staff to help with the increased workload.}
{Đây là một động từ dùng để chỉ việc thuê/tuyển dụng ai đó.}

\vocabentry{Job market}
{/ˈdʒɑːb ˈmɑːrkɪt/ (n)}
{The supply of people in a particular country or area who are able and willing to work in relation to the number of jobs available.}
{The job market is currently very difficult for recent graduates.}
{Đây là một cụm danh từ dùng để chỉ thị trường lao động, tình hình việc làm}

\vocabentry{Labor}
{/ˈleɪbər/ (n)}
{Work, especially hard physical work.}
{The cost of labor is a significant factor in the price of the final product.}
{Đây là một danh từ dùng để chỉ lao động hoặc nhân công.}

\vocabentry{Marketing}
{/ˈmɑːrkɪtɪŋ/ (n)}
{The action or business of promoting and selling products or services.}
{She has a successful career in marketing and communications.}
{Một lĩnh vực nghề nghiệp phổ biến.}

\vocabentry{Notice}
{/ˈnoʊtɪs/ (n)}
{Advance notification of intention to end an agreement or leave a job.}
{You must give at least two weeks' notice if you plan to resign.}
{Đây là một danh từ dùng để chỉ việc thông báo nghỉ việc.}

\vocabentry{Permanent}
{/ˈpɜːrmənənt/ (adj)}
{Lasting or intended to last or remain unchanged indefinitely.}
{After a successful probation period, she was offered a permanent position.}
{Thường dùng trong hợp đồng dài hạn - permanent contract.}

\vocabentry{Quality control}
{/ˈkwɑːləti kənˌtroʊl/ (n)}
{A system of maintaining standards in manufactured products by testing a sample of the output against the specification.}
{Every product must pass through quality control before it is shipped to customers.}
{Một lĩnh vực công việc quan trọng.}

\vocabentry{Retail}
{/ˈriːteɪl/ (n)}
{The sale of goods to the public in relatively small quantities for use or consumption rather than for resale.}
{She has worked in the retail industry for over ten years.}
{Một lĩnh vực nghề nghiệp.}

\vocabentry{Target}
{/ˈtɑːrɡɪt/ (n)}
{A person, object, or place selected as the aim of an attack.}
{We are on track to meet our sales targets for this quarter.}
{Trong công việc là chỉ tiêu cần đạt.}

\vocabentry{Underpaid}
{/ˌʌndərˈpeɪd/ (adj)}
{Paid too little.}
{Many social workers feel that they are underpaid and undervalued.}
{Đây là một tính từ dùng để chỉ tình trạng được trả lương quá thấp so với năng lực hoặc công việc.}

\vocabentry{Absenteeism}
{/ˌæbsənˈtiːɪzəm/ (n)}
{The practice of regularly staying away from work or school without good reason.}
{High levels of absenteeism can be a sign of low employee morale.}
{Đây là một danh từ dùng để chỉ tình trạng vắng mặt thường xuyên không lý do.}

\vocabentry{Blue-collar worker}
{/ˌbluː ˈkɑːlər ˈwɜːrkər/ (n)}
{A person who performs manual labor, especially in manufacturing or industrial work.}
{Many blue-collar workers in the construction industry are facing job losses.}
{Đây là một cụm danh từ dùng để chỉ công nhân lao động chân tay, đối lập với "white-collar worker" - nhân viên văn phòng}

\vocabentry{Contract}
{/ˈkɑːntrækt/ (n)}
{A written or spoken agreement, especially one concerning employment.}
{Please read the terms of your contract carefully before signing it.}
{Đây là một danh từ dùng để chỉ hợp đồng lao động.}

\vocabentry{Diversity}
{/daɪˈvɜːrsəti/ (n)}
{The state of being diverse; variety.}
{Promoting diversity and inclusion is a key goal for our HR department.}
{Trong công việc là sự đa dạng nhân sự.}

\vocabentry{Experience}
{/ɪkˈspɪriəns/ (n)}
{Practical contact with and observation of facts or events.}
{You need several years of experience in a similar role to apply for this position.}
{Đây là một danh từ dùng để chỉ kinh nghiệm làm việc.}

\vocabentry{Feedback}
{/ˈfiːdbæk/ (n)}
{Information about reactions to a product, a person's performance of a task, etc.}
{Regular feedback from your manager can help you improve your performance.}
{Đây là một danh từ dùng để chỉ phản hồi về công việc.}

\vocabentry{Incentive scheme}
{/ɪnˈsentɪv skiːm/ (n)}
{A plan or system for encouraging people to do something.}
{The company has introduced a new incentive scheme to boost sales.}
{Đây là một cụm danh từ dùng để chỉ hệ thống khen thưởng.}

\vocabentry{Job application}
{/ˈtʒɑːb ˌæplɪˈkeɪʃn/ (n)}
{A formal request to an authority for something.}
{You should tailor each job application to the specific requirements of the role.}
{Đây là một cụm danh từ dùng để chỉ việc nộp đơn xin việc.}

\vocabentry{Lay off}
{/ˌleɪ ˈɔːf/ (v)}
{Give up or stop doing something.}
{The company had to lay off several employees due to a decline in sales.}
{Trong công việc là cho nghỉ việc do công ty gặp khó khăn.}

\vocabentry{Mentee}
{/menˈtiː/ (n)}
{A person who is advised, trained, or counseled by a mentor.}
{As a mentee, you should be proactive in asking for advice and feedback.}
{Đây là một danh từ dùng để chỉ người được cố vấn.}

\vocabentry{Office}
{/ˈɔːfɪs/ (n)}
{A room, set of rooms, or building used as a place for commercial, professional, or bureaucratic work.}
{Our office is located in the heart of the city's financial district.}
{Đây là một danh từ dùng để chỉ văn phòng.}

\vocabentry{Resume}
{/rɪˈzuːmeɪ/ (n)}
{A brief account of a person's education, qualifications, and previous experience.}
{Make sure your resume is clear, concise, and easy to read.}
{Từ đồng nghĩa với CV, thường dùng trong tiếng Anh Mỹ.}

\vocabentry{Staff}
{/stæf/ (n)}
{All the people employed by a particular organization.}
{Our staff are our most valuable asset, and we invest heavily in their training.}
{Đây là một danh từ dùng để chỉ toàn bộ nhân viên.}

\vocabentry{Training}
{/ˈtreɪnɪŋ/ (n)}
{The action of teaching a person or animal a particular skill or type of behavior.}
{All new employees receive thorough training before they start their roles.}
{Đây là một danh từ dùng để chỉ hoạt động đào tạo.}

\vocabentry{Unemployed}
{/ˌʌnɪmˈplɔɪd/ (adj)}
{(Of a person) without a paid job but available to work.}
{He has been unemployed for six months and is struggling to find a new job.}
{Đây là một tính từ dùng để chỉ người thất nghiệp.}

\vocabentry{Voluntary work}
{/ˈvɑːlənteri wɜːrk/ (n)}
{Work that is done by choice and without being paid.}
{Doing voluntary work can be a great way to gain new skills and experience.}
{Đây là một cụm danh từ dùng để chỉ công việc tình nguyện.}

\vocabentry{Workload}
{/ˈwɜːrkloʊd/ (n)}
{The amount of work to be done by someone or something.}
{If you feel that your workload is too heavy, you should talk to your manager.}
{Đây là một danh từ dùng để chỉ khối lượng công việc.}

